\scanpage{90}
\mspage{89}
alors pour $\forall$ comp. connexe $T_i$ de $T$,
on a \deleted{$\pi_1(P_i,x_i)$}
\[
(5)\qquad \deleted{\unsure{1}}\quad
\pi\longrightarrow\pi_1(P_i,x_i)\mathbin{\text{---}}\pi_1(T_i,t_i)\mathbin{\text{---}}1
\]
\margin{$x_i\in P_i=P|T_i$\quad $t_i$ image dans $T_i$}
$\pi=\pi_1(U,1)$, supposons que les
hom $\pi\longrightarrow\pi_1(P_i,x_i)$ soient injectifs, alors ---
\[
\Pi_1(p):\quad \Pi_1(P)\mathbin{\text{---}}\Pi_1(T)
\]
\margin{$p:P\mathbin{\text{---}}T$\quad projection}
est muni d'une $\pi$-structure. (Si
$\pi\longrightarrow\pi_1(P_i,x_i)$ \deleted{est} n'est pas supposé injectif
et que son noyau $N_i$ est ind. de $i$,
alors $\Pi_1(p)$ est \deleted{muni d'} une $\pi'$-fibration,
\[
\pi'=\pi/N_i.\quad)
\]
\margin{\uline{NB} Si on remplace \unsure{le groupe constant}
$U$ par un \unsure{fibré loc. trivial} en groupes $U_T$ sur $T$
(défini par exemple par une opération de \unsure{$\Pi_1(T)$} sur
\unsure{un groupe connexe $U_0$}), \unsure{on a un groupe}
\unsure{local.} $\unsure{\underline{\pi}_t|\pi_t}\unsure{\simeq|=}\pi_1(U_t)$ \unsure{sur $T$}, et une
$\unsure{\pi_0|\underline{\pi}}$-structure sur $\Pi_1(P)$.}
\margin{\unsure{ici on peut} \deleted{\unsure{prendre $\mu$}}
\inserted{\unsure{remplacer $\mu$ par un}} \unsure{fibré en groupes} (6)
\unsure{sur $T$, et $E$ par un} \unsure{ensemble fibré} \unsure{relativement}.}

Soit, dans l'exemple précédent, $E$ un ensemble,
\inserted{$\mu$} un groupe \deleted{commutatif} qui opère sur
$E$,
\margin{\unsure{discret}\quad \deleted{\unsure{fibré}}}
\[
(6)\qquad \mu\xrightarrow{i_\mu}U
\]
un hom de groupes, et
\[
(7)\qquad \unsure{E}\xrightarrow{i_E}P
\]
une \deleted{hom} application, compatible avec les
opérations de $\mu$. Si on regarde le \deleted{\unsure{équi}} groupoïde
\inserted{$C''$} sur $E$ comme ens. d'objets, induit par
$\Pi_1(P)$ et (7),
\margin{\uline{NB} comme $\mu$ opère librement sur $P$,
(\unsure{quand} $i_\mu$ est injectif) il faut que $\mu$ opère
librement sur $E$.}
